% =============================================================================
% 题目与标题分页配置（高频设置）
% =============================================================================
% Book 模式在 Everflow-Book.tex 顶部切换：
%   standard / compact / loose / single / padl / padp
% Public 模式在 Everflow-Public.tex 顶部切换：
%   standard / compact / loose / single
%
% compact / standard / loose 只改变两题之间的空白；题干、选项和图片内部
% 间距保持一致。single、padl、padp 都是一题一页。

% ============================================================
% 【两题之间的空白】
%
% 数值越大，两道题之间的留白越多；数值越小，排版越紧凑。
% 所有长度必须写明单位，例如 cm、mm、pt 或 \baselineskip。
%
% 示例 1：把 standard 改得更紧凑。
% \def\EverflowQuestionGapStandard{1.2cm}
%
% 示例 2：让 Exam 中相邻题目稍微分开。
% \def\EverflowExamQuestionGap{0.7\baselineskip}
% ============================================================
\providecommand{\EverflowQuestionGapCompact}{0.35\baselineskip}
\providecommand{\EverflowQuestionGapStandard}{2.5cm}
\providecommand{\EverflowQuestionGapLoose}{7.5cm}
\providecommand{\EverflowExamQuestionGap}{0.45\baselineskip}

% ============================================================
% 【题目与选项是否尽量保持在同一页或同一栏】
%
% true：先把整道 bbox 题目装箱测量；放得下时整题移动，避免题干与选项拆开。
% false：允许题目自然分页，适合超过一整页的超长题目。
%
% OversizeQuestionPolicy 当前推荐保持 warn-break：
% 超高题目会给出警告，并自动改用可分页输出，避免形成超出版心的死盒。
% ============================================================
\providecommand{\EverflowKeepQuestionWithChoices}{true}
\providecommand{\EverflowOversizeQuestionPolicy}{warn-break}

% ============================================================
% 【题号与选项基础尺寸】
%
% 这些参数通常不需要修改。宽度增大后，题号或选项标签占用的横向空间增大，
% 正文可用宽度会相应减小。
%
% 错误写法：\def\EverflowQuestionLabelWidth{2.5}
% 原因：长度缺少 em、mm、pt 等单位。
% ============================================================
\providecommand{\EverflowQuestionLabelWidth}{2.25em}
\providecommand{\EverflowQuestionLabelSep}{0.55em}
\providecommand{\EverflowQuestionLeftIndent}{0pt}
\providecommand{\EverflowQuestionRightIndent}{0pt}
\providecommand{\EverflowChoiceLabelWidth}{2.25em}
\providecommand{\EverflowChoiceColumnGap}{1.2em}
\providecommand{\EverflowChoiceFourColumnThreshold}{0.22}
\providecommand{\EverflowChoiceTwoColumnThreshold}{0.46}
\providecommand{\EverflowChoiceBeforeSkip}{0.45\baselineskip}
\providecommand{\EverflowChoiceRowGap}{0.22em}

% =============================================================================
% 标题强制换页与“标题不能孤立”保护
% =============================================================================
% 这两类功能彼此独立：
%
% ForceNewPage=true：该标题在开始一个新标题组时强制进入下一张物理页。
% KeepWithNext=true：不强制换页，但标题不能单独留在页尾或栏尾。
%
% 推荐组合：
% \def\EverflowBookSectionForceNewPage{true}
% \def\EverflowExamSectionForceNewPage{false}
% \def\EverflowPublicSectionForceNewPage{true}
% \def\EverflowSectionKeepWithNext{true}
%
% 结果：Book 与 Public 的 section 默认从下一张物理页开始；Exam 保持现有
% 连续排版。当前页或当前栏空间不足时，标题组仍会与后面的第一项有效内容一起移动。

% ============================================================
% 【全部标题防孤立总开关】
%
% true：启用连续标题组与首个有效内容的联合分页。
% false：关闭全部智能跟随；各标题的 KeepWithNext 设置暂时不生效。
%
% 示例 1：正常启用。
% \def\EverflowHeadingKeepWithNextControl{true}
%
% 示例 2：排查第三方宏包冲突时临时关闭。
% \def\EverflowHeadingKeepWithNextControl{false}
%
% true 和 false 必须使用小写英文，不能写 True、TRUE、是或否。
% ============================================================
\providecommand{\EverflowHeadingKeepWithNextControl}{true}

% ============================================================
% 【连续标题是否合并为一个标题组】
%
% true：questiongroup、chapter、section、subsection、题型说明等连续标题只进行
% 一次分页判断，不会被拆到不同页面或不同栏位。
%
% false：每个标题独立判断。仅用于兼容极少数特殊旧题源。
% ============================================================
\providecommand{\EverflowHeadingConsecutiveGroupControl}{true}

% ============================================================
% 【标题后至少保护多少行普通正文】
%
% 当标题后不是 bbox、图片、表格或代码块，而是普通正文时使用该参数。
%
% 3：至少保护正文开头约 3 行，推荐值。
% 5：保护更强，但页尾留白可能增加。
%
% 合法建议范围为 1--5。不要写 3行，也不要添加长度单位。
% ============================================================
\providecommand{\EverflowHeadingMinimumFollowLines}{3}

% ============================================================
% 【标题后优先保护的内容模式】
%
% question：优先保护第一道 bbox 题目；适合试卷、题库和习题册，推荐使用。
% block：把第一项图片、表格、代码或公式块作为完整内容块保护。
% lines：普通正文至少保护 HeadingMinimumFollowLines 指定的行数。
%
% 当前统一内核会自动识别这三类内容；该参数用于表达默认偏好并保留扩展接口。
% ============================================================
\providecommand{\EverflowHeadingFollowMode}{question}

% ============================================================
% 【Exam A3 双栏中标题组空间不足时前进到哪里】
%
% column：左栏不足时先进入右栏；右栏不足时进入下一张 A3 的左栏。
% 这是页面利用率最高的默认方式。
%
% page：当前栏的预留空间不足时直接进入下一张 A3 物理页。
% 适合要求大标题尽量从新纸张开始的版式。
%
% 错误写法：\def\EverflowExamAThreeHeadingAdvance{right}
% 原因：合法值只有 column 和 page。
% ============================================================
\providecommand{\EverflowExamAThreeHeadingAdvance}{column}

% ============================================================
% 【各种标题是否强制另起物理页】
%
% true：标题在开始一个新标题组时执行物理换页。
% false：允许连续排版，再由 KeepWithNext 判断是否需要移动。
%
% section 使用三版独立默认值：Book/Public=true，Exam=false。
% 旧全局参数 EverflowSectionForceNewPage 默认写 mode，表示服从三版独立值；
% 若旧工程在主入口预设 true 或 false，则仍可一次覆盖三版，保证兼容。
% 其他标题默认 false，避免连续标题组中途重复换页。
% ============================================================
\providecommand{\EverflowQuestionGroupForceNewPage}{false}
\providecommand{\EverflowPartForceNewPage}{false}
\providecommand{\EverflowChapterForceNewPage}{false}
\providecommand{\EverflowBookSectionForceNewPage}{true}
\providecommand{\EverflowExamSectionForceNewPage}{false}
\providecommand{\EverflowPublicSectionForceNewPage}{true}
\providecommand{\EverflowSectionForceNewPage}{mode}
\providecommand{\EverflowSubsectionForceNewPage}{false}
\providecommand{\EverflowSubsubsectionForceNewPage}{false}
\providecommand{\EverflowPublicationUnitTitleForceNewPage}{false}
\providecommand{\EverflowExamSectionNoteForceNewPage}{false}

% ============================================================
% 【各种标题是否与后面的第一项有效内容一起分页】
%
% true：该标题参与连续标题组，不能单独留在 A4 页尾、Exam 页尾、A3 左栏尾
% 或 A3 右栏尾。
%
% false：该标题不主动建立保护链，可能按自然分页停在页尾或栏尾。
%
% 示例 1：只关闭 subsubsection 的保护。
% \def\EverflowSubsubsectionKeepWithNext{false}
%
% 示例 2：全部保持推荐值。
% 下列八项全部使用 true。
% ============================================================
% 新名称 KeepWithNext 的优先级最高；如果旧题源只预设了 OrphanControl
% 或 AutoPageBreak，则在这里把旧值复制到新名称，保证旧入口仍然有效。
% 这一段是兼容层，不需要手工修改。
\newcommand{\EverflowResolveLegacyHeadingSwitch}[3]{%
  \ifcsname #1\endcsname
  \else
    \ifcsname #2\endcsname
      \expandafter\edef\csname #1\endcsname{\csname #2\endcsname}%
    \else
      \ifcsname #3\endcsname
        \expandafter\edef\csname #1\endcsname{\csname #3\endcsname}%
      \else
        \expandafter\def\csname #1\endcsname{true}%
      \fi
    \fi
  \fi
}
\EverflowResolveLegacyHeadingSwitch{EverflowQuestionGroupKeepWithNext}{EverflowQuestionGroupOrphanControl}{EverflowQuestionGroupAutoPageBreak}
\EverflowResolveLegacyHeadingSwitch{EverflowPartKeepWithNext}{EverflowPartOrphanControl}{EverflowPartAutoPageBreak}
\EverflowResolveLegacyHeadingSwitch{EverflowChapterKeepWithNext}{EverflowChapterOrphanControl}{EverflowChapterAutoPageBreak}
\EverflowResolveLegacyHeadingSwitch{EverflowSectionKeepWithNext}{EverflowSectionOrphanControl}{EverflowSectionAutoPageBreak}
\EverflowResolveLegacyHeadingSwitch{EverflowSubsectionKeepWithNext}{EverflowSubsectionOrphanControl}{EverflowSubsectionAutoPageBreak}
\EverflowResolveLegacyHeadingSwitch{EverflowSubsubsectionKeepWithNext}{EverflowSubsubsectionOrphanControl}{EverflowSubsubsectionAutoPageBreak}
\EverflowResolveLegacyHeadingSwitch{EverflowPublicationUnitTitleKeepWithNext}{EverflowPublicationUnitTitleOrphanControl}{EverflowPublicationUnitTitleAutoPageBreak}
\EverflowResolveLegacyHeadingSwitch{EverflowExamSectionNoteKeepWithNext}{EverflowExamSectionNoteOrphanControl}{EverflowExamSectionNoteAutoPageBreak}

% ============================================================
% 【兼容旧版 AutoPageBreak 与 OrphanControl 名称】
%
% 旧工程在主入口加载前写入旧名称时仍会生效。新工程应优先使用 KeepWithNext，
% 因为新名称更准确地表达“标题与后续内容保持在一起”的含义。
% ============================================================
\providecommand{\EverflowQuestionGroupAutoPageBreak}{\EverflowQuestionGroupKeepWithNext}
\providecommand{\EverflowQuestionGroupOrphanControl}{\EverflowQuestionGroupKeepWithNext}
\providecommand{\EverflowPartAutoPageBreak}{\EverflowPartKeepWithNext}
\providecommand{\EverflowPartOrphanControl}{\EverflowPartKeepWithNext}
\providecommand{\EverflowChapterAutoPageBreak}{\EverflowChapterKeepWithNext}
\providecommand{\EverflowChapterOrphanControl}{\EverflowChapterKeepWithNext}
\providecommand{\EverflowSectionAutoPageBreak}{\EverflowSectionKeepWithNext}
\providecommand{\EverflowSectionOrphanControl}{\EverflowSectionKeepWithNext}
\providecommand{\EverflowSubsectionAutoPageBreak}{\EverflowSubsectionKeepWithNext}
\providecommand{\EverflowSubsectionOrphanControl}{\EverflowSubsectionKeepWithNext}
\providecommand{\EverflowSubsubsectionAutoPageBreak}{\EverflowSubsubsectionKeepWithNext}
\providecommand{\EverflowSubsubsectionOrphanControl}{\EverflowSubsubsectionKeepWithNext}
\providecommand{\EverflowPublicationUnitTitleAutoPageBreak}{\EverflowPublicationUnitTitleKeepWithNext}
\providecommand{\EverflowPublicationUnitTitleOrphanControl}{\EverflowPublicationUnitTitleKeepWithNext}
\providecommand{\EverflowExamSectionNoteAutoPageBreak}{\EverflowExamSectionNoteKeepWithNext}
\providecommand{\EverflowExamSectionNoteOrphanControl}{\EverflowExamSectionNoteKeepWithNext}

% ============================================================
% 【标题组开始时预留的基础空间】
%
% 这些数值是标题组尚未测量到第一项内容时使用的安全预留量。
% 第一项 bbox 会在完成装箱后使用真实高度继续保护。
%
% 数值越大，防孤立越稳健，但页尾留白可能增加；数值越小，页面更紧凑，
% 但非常高的自定义标题可能需要单独增大对应参数。
% ============================================================
\providecommand{\EverflowQuestionGroupNeedSpace}{10\baselineskip}
\providecommand{\EverflowPartNeedSpace}{10\baselineskip}
\providecommand{\EverflowChapterNeedSpace}{8\baselineskip}
\providecommand{\EverflowSectionNeedSpace}{9\baselineskip}
\providecommand{\EverflowSubsectionNeedSpace}{6\baselineskip}
\providecommand{\EverflowSubsubsectionNeedSpace}{5\baselineskip}
\providecommand{\EverflowPublicationUnitTitleNeedSpace}{9\baselineskip}
\providecommand{\EverflowSectionNoteNeedSpace}{4\baselineskip}

% =============================================================================
% Exam 科目标题与第一个 section 的显示关系
% =============================================================================
% inline：科目标题和第一个 section 连续出现，并进入同一个标题组。
% standalone：科目标题独占一页，第一个 section 再按自己的设置处理。
% hidden：Exam 不显示科目标题，但仍写入目录。
%
% Merge=true 时，questiongroup 与紧随其后的第一个 section 只建立一个标题组，
% 不会重复执行物理换页。
% =============================================================================
\providecommand{\EverflowExamQuestionGroupMode}{inline}
\providecommand{\EverflowExamMergeQuestionGroupWithFirstSection}{true}

% Exam Part/科目标题图片开关，只控制 questiongroup 的第三个图片参数。
% false：隐藏 \questiongroup{数据结构}[图片路径][宽度] 中的装饰图片，标题仍显示。
% true：显示该装饰图片。
% 题干、选项、答案解析、普通 includegraphics、封面、Logo、二维码和水印均不受影响。
\providecommand{\EverflowExamPartImageEnabled}{false}

% 旧名称继续兼容；新工程优先修改 EverflowExamPartImageEnabled。
\providecommand{\EverflowExamQuestionGroupShowImage}{\EverflowExamPartImageEnabled}
